\chapter{Structure and Format Text}
\label{chap:text}

\begin{learningoutcomes}
  \item Build a clearly structured article.
  \item Use emphasis, lists, quotations, and footnotes.
  \item Define a command once and reuse it.
\end{learningoutcomes}

\section{Begin with structure, not appearance}

This complete file covers the text tools used in most reports:

\begin{latexcode}{Complete example: structured text}
\documentclass{article}

\title{Growth Under Two Conditions}
\author{Your Name}
\date{}

\begin{document}
\maketitle

\begin{abstract}
We compare growth under two simulated conditions.
\end{abstract}

\section{Introduction}
The aim is to compare \emph{Condition A} with
\emph{Condition B}.

\section{Materials and Methods}
The work followed three steps:
\begin{enumerate}
  \item prepare the samples;
  \item record the measurements; and
  \item compare the group means.
\end{enumerate}

\section{Results}
The main observations were:
\begin{itemize}
  \item growth increased in both groups;
  \item Condition B had the larger mean; and
  \item no values were missing.
\end{itemize}

\section{Discussion}
\textbf{Main result.} Condition B showed greater growth.
The result is descriptive, not proof of causation.

\section{Conclusion}
The example demonstrates a compact scientific structure.
\end{document}
\end{latexcode}

Use section commands to express hierarchy. Use \cmd{emph} for contextual
emphasis and \cmd{textbf} for a short label. Do not format every important
sentence manually.

\section{Choose the right list}

\begin{latexcode}{List recipes}
% Order matters:
\begin{enumerate}
  \item Collect data.
  \item Check data.
  \item Analyse data.
\end{enumerate}

% Order does not matter:
\begin{itemize}
  \item source file
  \item figure file
  \item bibliography file
\end{itemize}

% Term and explanation:
\begin{description}
  \item[Input] The source supplied to a process.
  \item[Output] The result produced by the process.
\end{description}
\end{latexcode}

\section{Quotations and notes}

\begin{latexcode}{Quotation and footnote}
Short quoted words remain inside the paragraph.

\begin{quote}
A longer quotation is set apart from the surrounding text.
\end{quote}

A footnote holds a brief supporting note.\footnote{Keep
essential reasoning in the main text.}
\end{latexcode}

\section{Define repeated material once}

Put custom commands in the preamble:

\begin{latexcode}{Complete example: reusable commands}
\documentclass{article}

\newcommand{\species}[1]{\textit{#1}}
\newcommand{\samplecount}{24}
\newcommand{\studyname}{Growth Trial}

\begin{document}
\section{\studyname}

We measured \samplecount{} samples of
\species{Arabidopsis thaliana}.

The \studyname{} used the same sample count in both groups.
\end{document}
\end{latexcode}

When a value changes, edit its definition once. Give personal commands
distinctive names so they do not collide with existing commands.

\section{Control page breaks sparingly}

\begin{latexcode}{Useful page commands}
\newpage      % finish the current page
\clearpage    % finish the page and place waiting figures/tables
\end{latexcode}

Let \LaTeX{} choose ordinary line and page breaks. Forced breaks are for
deliberate structural boundaries, not routine spacing.

\begin{codelab}
Start with the structured-text example. Replace the topic, add a description
list, and define one command for a value repeated three times. Compile after
each change.
\end{codelab}

\begin{chapterchecklist}
  \item Headings describe the argument rather than its appearance.
  \item Lists are used only when they make scanning easier.
  \item Repeated terms or values are defined once.
  \item Manual spacing and forced breaks are exceptional.
\end{chapterchecklist}
